\documentclass[11pt,a4paper]{article}
\usepackage[margin=25mm]{geometry}
\usepackage{amsmath,amssymb,amsthm,hyperref}
\setlength{\emergencystretch}{2em}
\newtheorem{lemma}{Lemma}\newtheorem{theorem}[lemma]{Theorem}
\newcommand{\gen}[1]{\langle#1\rangle}
\newcommand{\Sym}{\operatorname{Sym}}
\title{The two ends of the critical configuration}
\author{Human Compression project: a component of the L019 recovery}
\date{}
\begin{document}\maketitle
We retain the standing hypotheses and strict smaller-degree induction
of the acquired first-cut component. Thus the full fractured natural
symmetric string C-marking consists of an even degree-six rank-three
primitive $A_5$ factor $E$ on $Y$, followed by
\[
 P=\gen{t_0,\ldots,t_{s-1}}=\Sym(X),\qquad
 |X|=m,\quad s=m/2+1,\quad X\cap Y=\{a\}.
\]
The common point is an endpoint anchor in each factor; $m$ is even.
The full degree is $n=m+5$ and its rank is $m/2+4$.
A sufficient split has the precise meaning fixed in the acquired
component: it is pure, or has a trivial internal factor and two odd
actual one-point attaching groups. The induction applies only to
smaller degrees satisfying the strict inequality $2u>d+2$.

We may assume the least unique crossing of $P$ has index zero;
the acquired first-cut theorem has disposed of every positive index.
Write
\[
 t_0=(a\ b)\eta,\qquad T=\gen{t_1,\ldots,t_{s-1}}.
\]
Its perfect split has singleton side $\{a\}$ and complementary
side $X\setminus\{a\}$. The actual-wing bound, since
$2\operatorname{rank}(T)=m>m-1$, gives
$T=\Sym(X\setminus\{a\})$.

\begin{lemma}\label{tail}
Either the full marking has an internal sufficient split, or $T$
has a least unique crossing at an index $k\ge2$ in its own order.
It is a perfect split of $T$ and inherits to the full marking
at the original $P$ index $j=k+1$. If $A',B$ are its two
$T$-omission sides, with $b\in A'$, then
\[
 |A'|\ge 2k+1,\quad
 \eta\text{ is supported on }A',\quad
 a_0:=|\{a\}\cup A'|\ge 2j.
                                                        \tag{1}
\]
If $a_0>2j$, the full marking has a sufficient split.
\end{lemma}
\begin{proof}
The tail is fractured, inherits IP as an actual supported subgroup,
and has endpoint anchor $b$. Its degree is $m-1$ and twice its
rank is $m$, so the acquired anchored first-split theorem applies.
It supplies the perfect split and $2k\le |A'|-1$.
If $k=0$, the adjacent inheritance lemma applies with tail rank
$m/2$ and degree $m-1$: its strict condition is $m>m-2$.
It forces $\eta=1$. The original mark $t_0$ is then pure and
is internal in the full marking, because three $E$ marks precede
it and the nonempty tail follows it. This is sufficient.
If $k=1$, the one-mark left wing of $T$ is transitive of
degree two, contrary to $|A'|\ge3$.

For $k\ge2$, nonadjacent inheritance at the original index zero
applies to $T$'s split. In its actual support proof, $t_0$ fixes
$B$ pointwise, so $\eta$ is supported on $A'$. The new left
side of $P$ is $\{a\}\cup A'$, giving (1).
The same split in the full marking has sides $Y\cup A'$ and
$B$: the $E$ factor attaches only at $a$, and the split mark
fixes $Y$.

Put $\delta=a_0-2j$. When $\delta\ge1$, the acquired
positive-defect proof applies verbatim, although $j$ need not
be the least unique index of $P$. To make this interface explicit,
that proof after its first lemma uses only a perfect split, $j\ge2$,
and $a_0\ge2j+1$. These give right-wing rank $q$ with
$2q=|B|+\delta>|B|$, hence an actual symmetric right wing.
The remaining faithful-augmentation and parity cases use exactly
these quantities and the full strict parent. They therefore provide
the same sufficient split under the same smaller-degree induction.
\end{proof}

\begin{theorem}\label{ends}
If the full marking has no internal sufficient split, then for
some $k\ge2$ its labels and actual supports have the form
\[
                       E,\ q_-,\ h_0,\ldots,h_{k-1},\ q_+,\ E'.
                                                               \tag{2}
\]
There are pairwise disjoint sets $Y,W,Z$, of sizes $6,2k+1,6$,
on which $E,M:=\gen{h_0,\ldots,h_{k-1}},E'$ are respectively
supported and transitive. Both $E,E'$ are even primitive
degree-six rank-three $A_5$ factors. For $a\in Y$, $d\in Z$
and $b,c\in W$,
\[
 q_-=(a\ b)\eta,\qquad q_+=(c\ d)\alpha,
                                                               \tag{3}
\]
where $\eta,\alpha$ are nonidentity involutions supported on
$W$, fixing $b,c$ respectively. The domain has size $2k+13$
and the rank is $k+8$. Moreover, as actual supported groups,
\[
\begin{aligned}
 M&=\Sym(W),\\
 \gen{q_-,M}&=\Sym(W\cup\{a\}),&
 \gen{M,q_+}&=\Sym(W\cup\{d\}),\\
 \gen{q_-,M,q_+}&=\Sym(W\cup\{a,d\}).
\end{aligned}                                                  \tag{4}
\]
The marking of $M$ is fractured, has endpoint anchors $b,c$,
and has at least two eligible edges of every colour.
\end{theorem}
\begin{proof}
By Lemma~\ref{tail}, only $\delta=0$ remains. Put $W=A'$,
so $|W|=2k+1$, $j=k+1$, and $a_0=2j$.
Set $h_i=t_{i+1}$ for $0\le i<k$, and $q_-=t_0$, $q_+=t_j$.
The left full wing
\[
 F=\gen{E,t_0,\ldots,t_{j-1}}
\]
has degree $f=a_0+5$ and rank $j+3$, with $2(j+3)=f+1$.
The actual-wing bound forces $F=\Sym(Y\cup W)$.
Write the right wing as $R$, with degree $z=|B|$ and rank $q$.
The critical rank identity gives $2q=z$. In particular $z$ is
positive and even. Write the split mark
\[
                  t_j=(c\ d)\alpha\beta
\]
with internal factors on the two full sides. Its left factor
fixes $Y$, since $t_j$ does, and is supported on $W$.

If both factors are nonidentity, faithful augmentation with
symmetric wing $F$ gives a proper primitive faithful C-marking
of degree $z+1$ and rank $q+1$. Its rank is strictly above half
its degree. The acquired primitive classification allows proper
exceptions only in degrees six and ten, whereas $z+1$ is odd.
This is impossible, including the small odd degrees; natural
alternating and symmetric groups are excluded by the same
augmentation theorem.

Suppose $\alpha=1$. The attaching groups are now $F$ and
$J=\gen{t_j,R}$, meeting at $c$. Since $F$ is odd, an odd
$J$ gives a sufficient split. If $J$ is even, the acquired
even-attachment theorem gives either $2\operatorname{rank}(F)>f+2$
or the exceptional right degree $z=5$. The former contradicts
$2\operatorname{rank}(F)=f+1$ and the latter contradicts the
evenness of $z$. This case is disposed of.

Consequently $\beta=1$. The group $H=\gen{F,t_j}$ is
$\Sym(Y\cup W\cup\{d\})$: its symmetric subgroup $F$
contains $\alpha$, so it also contains the single crossing
transposition and hence generates that full symmetric group.
It meets $R$ only at $d$. If $R$ is odd, this is a sufficient
split. Otherwise $R$ is even, fractured, anchored at $d$, and
has rank $z/2$. The acquired anchored-even theorem excludes
$z\ge12$, and the small-even theorem applies for $z=6,8,10$
with the actual odd external group $H$. Degrees two and four
cannot be even transitive C-markings of the required ranks:
degree two has no even involution, while every nonidentity even involution of degree four is a
fixed-point-free double transposition, incompatible with the anchor
fixed by the second mark. Thus the only
remaining case is $z=6$, $q=3$, and $R=E'$ is the stated
primitive $A_5$. This proves (2)--(3). Nonidentity of both
remaining internal factors follows because otherwise one of
the two connecting marks would already be a pure internal split.

It remains to identify the middle groups without projecting
any C-marking. The original $P$ and $T$ are the actual groups
\[
 P=\Sym(\{a\}\cup W\cup Z),\qquad T=\Sym(W\cup Z).
\]
We have just proved
$F=\Sym(Y\cup W)$ and $H=\Sym(Y\cup W\cup\{d\})$.
The intersections of supported symmetric groups are the symmetric
groups on their set intersections. The original IP, applied to
their exact label sets, therefore gives
\[
\begin{array}{ll}
 F\cap T=M=\Sym(W),&
 F\cap P=\gen{q_-,M}=\Sym(W\cup\{a\}),\\
 H\cap T=\gen{M,q_+}=\Sym(W\cup\{d\}),&
 H\cap P=\gen{q_-,M,q_+}=\Sym(W\cup\{a,d\}).
\end{array}
\]
These are literal subgroup intersections in the given ambient
group; no quotient-IP assertion or new finite exception test enters.

The unique crossings in (3) and the string relations show that
all middle marks except $h_0$ fix $b$, and all except $h_{k-1}$
fix $c$. Their transitivity makes the respective end marks move
these points. Thus $b,c$ are endpoint anchors. They are distinct
because $k\ge2$: otherwise every middle mark would fix their
common point. Fracture descends from $T$, since a transitive
middle omission, together with $q_+$ and $E'$, would give a
transitive omission of $T$.

Finally $k$ was the least unique crossing index in $T$.
Every preceding colour therefore has at least two eligible
edges in $T$. Every edge of those marks lies in $W$.
Their eligibility persists upon passing to the smaller omission
subgroup of $M$. This proves the final assertion and completes
the reduction.
\end{proof}

\paragraph{Attribution and remaining task.}
The source extension problem and the acquired classifications remain
attributed to Cameron, Fernandes, Leemans and Mixer through their
individual recovered statements. This deduction uses the new first-cut
density and actual intersections to isolate the two-ended configuration.
In particular, (4) avoids a further primitive classification or finite
small-degree calculation. The remaining middle marking is symmetric
of degree $2k+1$ and rank $k$, with two endpoint anchors and paired
eligible edges. It has not been classified or excluded here. Neither
the full extension assertion nor its inverse counting correspondence
is claimed. There is no whole-original, conjecture, Lean or measurement
credit in this component.
\end{document}
